%%=============================================================================
%% Technologische context en keuze
%%=============================================================================

\chapter{\IfLanguageName{dutch}{Technologie en verantwoording}{Technology and rationale}}%
\label{ch:technologie}

%% Richtlengte: 2 tot 3 bladzijden.
%%
%% LET OP: DIT IS GEEN LITERATUURSTUDIE.
%% Je hoeft geen overzicht te schrijven van wat er in het vakgebied bestaat,
%% en je hoeft geen wetenschappelijke bronnen te verzamelen. De vraag hier is
%% praktisch: welke technologie heb je gebruikt, welke alternatieven waren er,
%% en waarom is het deze geworden?
%%
%% Dit is ook wat de jury je tijdens je verdediging zal vragen. Wie het hier
%% opschrijft, staat daar sterker.
%%
%% WEL: vermeld de bronnen die je echt gebruikt hebt. Voor een project zijn dat
%% meestal officiële documentatie, handleidingen, blogberichten van makers,
%% normen, of een boek. Vermeld ze omdat je lezer ze moet kunnen terugvinden,
%% niet omdat een lijst nu eenmaal lang moet zijn. Zie de uitleg onderaan dit
%% bestand en de voorbeelden in bronnen.bib.
TODO!
Het graduaatsproject werd uitgevoerd binnen een bestaande professionele softwareomgeving. Daardoor was de technologische vrijheid beperkter dan bij een volledig nieuw project. De belangrijkste technologieën en architecturale keuzes lagen reeds vast binnen Vesalius.ai. Dit was een bewuste keuze van het bedrijf: nieuwe functionaliteit moet aansluiten bij de bestaande applicatie, ontwikkelprocessen en infrastructuur.
De technologiekeuze binnen dit project moet daarom vooral worden bekeken vanuit onderhoudbaarheid, integratie met de bestaande applicatie, kennis binnen het team en aansluiting bij de bestaande architectuur.

\section{\IfLanguageName{dutch}{Bestaande omgeving}{Existing environment}}%
\label{sec:bestaande-omgeving}
TODO!
De backend van het Vesalius-platform is opgebouwd met PHP en Laravel. Laravel voorziet onder andere de structuur voor de backendlogica, routing en interactie met de database. Voor het graduaatsproject werd deze bestaande backendtechnologie behouden, omdat de nieuwe functionaliteit geïntegreerd moest worden in de bestaande applicatie.
De frontend is opgebouwd met Angular en TypeScript. HTML wordt gebruikt voor de structuur van de gebruikersinterface en SCSS voor de styling. Angular sluit aan bij de bestaande frontendarchitectuur en biedt een componentgebaseerde manier om nieuwe configuratiemogelijkheden en gebruikersinteracties toe te voegen.
Voor persistente gegevens wordt MySQL gebruikt. De configuratie van templates en de bijkomende gegevens die nodig zijn om de nieuwe functionaliteit te ondersteunen, moeten daarom aansluiten bij het bestaande datamodel en de bestaande databaseconventies.
Voor authenticatie en autorisatie wordt Keycloak gebruikt. Hierdoor hoeft het project geen eigen authenticatiesysteem te implementeren en kan de bestaande gebruikers- en toegangsstructuur worden gebruikt. Dit is vooral relevant voor functionaliteit die afhankelijk is van de ingelogde gebruiker, zoals het bepalen van de gekoppelde arts bij automatische generatie.
Redis wordt binnen de bestaande omgeving gebruikt voor caching en gerelateerde tijdelijke gegevens. Docker wordt gebruikt om onderdelen van de ontwikkelomgeving reproduceerbaar te maken.
De ontwikkeling gebeurt met Git en Bitbucket voor versiebeheer. Jira wordt gebruikt voor het opvolgen van tickets en werk. Tijdens de ontwikkeling werd gewerkt met Visual Studio Code en de bestaande ontwikkelomgeving, waaronder Ubuntu via WSL2 op Windows. De infrastructuur van het platform maakt daarnaast gebruik van Kubernetes, Terraform en CodeFresh voor respectievelijk containerorkestratie, infrastructuurbeheer en CI/CD.

\section{\IfLanguageName{dutch}{Vergelijking van de mogelijkheden}{Comparison of the options}}%
\label{sec:vergelijking}
TODO! (2-4 mogelijkheden vgl op criteria die voor dit project belangrijk zijn: aansluiting bestaande omgeving/leercurve/ondersteuning/licentie/prestaties/onderhoudbaarheid/kostprijs)
werk met tabel

Omdat het project binnen een bestaande professionele applicatie werd uitgevoerd, was een volledig nieuwe technologiekeuze niet aangewezen. Een nieuwe backendtechnologie of frontendframework zou de functionaliteit technisch wel kunnen realiseren, maar zou niet aansluiten bij de bestaande codebase.
De belangrijkste alternatieven kunnen daarom als volgt worden bekeken:

\begin{table}[htbp]
  \centering
  \begin{tabular}{lccc}
    \toprule
    \textbf{Criterium} & \textbf{Bestaande stack} & \textbf{Nieuwe stack} & \textbf{Losstaande oplossing} \\
    \midrule
    Aansluiting bestaande applicatie & ++ & -- & -- \\
    Onderhoudbaarheid door team & ++ & - & -- \\
    Integratie met bestaande data & ++ & - & -- \\
    Leercurve binnen projectperiode & ++ & - & + \\
    Hergebruik bestaande infrastructuur & ++ & -- & - \\
    \bottomrule
  \end{tabular}
  \caption[Vergelijking van technologische mogelijkheden]{Vergelijking van de bestaande technologische stack met alternatieve benaderingen.}
  \label{tab:vergelijking}
\end{table}

De bestaande stack scoort vooral beter omdat de functionaliteit rechtstreeks onderdeel wordt van het bestaande Vesalius-platform. Een aparte applicatie zou extra authenticatie, gegevensuitwisseling, deployment en onderhoud vereisen. Ook het gebruik van een nieuwe programmeertaal of framework zou weinig meerwaarde bieden voor een relatief beperkte uitbreiding van een bestaande applicatie.


\section{\IfLanguageName{dutch}{Vergelijking met de opleiding}{Comparison with the programme}}%
\label{sec:vergelijking-opleiding}

%% Verwacht onderdeel van je verdediging: hoe verhoudt de technologie die je
%% hier gebruikte zich tot wat je in de opleiding gezien hebt? Wat was
%% hetzelfde, wat was nieuw, en wat bleek in de praktijk anders te werken dan
%% in de les?
TODO+herschrijven! welke specifieke technieken? wat in praktijk anders (git versiecontrole via BitBucket/SourceTree bv en niet GitHub, planning gebeurt in JIRA (Atlassian) en bv niet in Github planning of Azure sprints, Google Chat voor te communiceren met elkaar ipv MsTeams, ...), wat nieuw (vanalles) en wat hetzelfde (git versiecontrole commands)? 
Het project sluit goed aan bij de competenties uit de opleiding Graduaat in het Programmeren. Tijdens de opleiding werd gewerkt met programmeerconcepten, databases, webontwikkeling, versiebeheer, softwarearchitectuur en samenwerking binnen softwareprojecten.
In de professionele omgeving werd dezelfde basiskennis toegepast binnen een bestaande en aanzienlijk grotere codebase. Het belangrijkste verschil met oefeningen en projecten binnen de opleiding is dat wijzigingen niet geïsoleerd bestaan. Elke aanpassing moet rekening houden met bestaande functionaliteit, gebruikers, gegevens, beveiliging, codeconventies en andere onderdelen van het platform.
Ook het werken met bestaande infrastructuur en bedrijfsprocessen vormde een verschil. Technologieën zoals Keycloak, Kubernetes, Terraform en CI/CD handling via CodeFresh  zijn volledig nieuwe zaken voor mij en worden niet enkel als afzonderlijke technologieën gebruikt, maar maken deel uit van een groter geheel waarin verschillende componenten samenwerken.


\section{\IfLanguageName{dutch}{Gemaakte keuze}{The choice made}}%
\label{sec:keuze}
TODO! 
Voor het graduaatsproject werd bewust gekozen om de bestaande Vesalius.ai-stack volledig te behouden. De backendfunctionaliteit werd geïntegreerd in Laravel/PHP, de gebruikersinterface werd uitgebreid met Angular en TypeScript en de bestaande MySQL-database werd gebruikt voor persistente gegevens. Keycloak bleef verantwoordelijk voor authenticatie en autorisatie.
Deze keuze vermijdt onnodige technische complexiteit en zorgt ervoor dat de functionaliteit aansluit bij de bestaande kennis van het team. Bovendien kan de oplossing gebruikmaken van bestaande infrastructuur, deploymentprocessen en ontwikkelafspraken.
Het belangrijkste nadeel van deze keuze is dat er minder ruimte is om nieuwe technologieën te introduceren. Voor dit project weegt die beperking echter niet op tegen de voordelen van een consistente integratie in het bestaande platform.

LEES GRONDIG: 
%% Eén heldere paragraaf: dit is het geworden, en dit is de reden. Vermeld
%% ook eerlijk wat je ervoor hebt moeten opgeven.
%% Verwijs naar je bronnen met \autocite{sleutel} of \textcite{sleutel}. De
%% sleutels staan in bronnen.bib. Zet de verwijzing binnen de zin, vlak bij de
%% bewering waar ze bij hoort - niet achteraan de paragraaf. Hieronder een
%% voorbeeld van beide vormen; vervang ze door je eigen bronnen.

Bij het beoordelen van leesbaarheid en onderhoudbaarheid van de code werd
uitgegaan van de gangbare vuistregels uit het vakgebied~\autocite{Martin2008}.
De officiële documentatie van het gekozen framework~\autocite{DockerDocs2026}
diende als voornaamste naslagwerk. \textcite{Fowler2018} gebruik je wanneer de
naam van de auteur zelf een rol speelt in je zin.

%%---------- Bronnen vermelden in een projectrapport ---------------------------
%%
%% Welke bronnen horen in je bronnenlijst?
%%   - officiële documentatie van de talen, frameworks en diensten die je gebruikt
%%   - handleidingen, blogberichten of video's waarop je je aanpak gebaseerd hebt
%%   - normen en standaarden waaraan je je gehouden hebt
%%   - een boek of artikel, als je er een gebruikt hebt
%%   - interne documenten van je bedrijf: vermeld ze, maar neem geen
%%     vertrouwelijke inhoud over zonder toestemming van je mentor
%%
%% Wat er NIET in hoort:
%%   - een gesprek met een AI-hulpmiddel. Dat is geen bron die iemand kan
%%     nagaan; het hoort thuis in de bijlage over je AI-gebruik.
%%   - bronnen die je niet gelezen hebt. Een lange lijst maakt geen indruk;
%%     een lijst met werken die nergens aangehaald worden, wel - de verkeerde.
%%
%% Vermeld bij een webbron ALTIJD de datum waarop je ze geraadpleegd hebt
%% (het veld urldate in bronnen.bib): documentatie verandert.
